% ============================================================================
%  E/19-bon-nhanh-luyen — file chương
%  Ngày tạo: 2026-09-07 | Sửa gần nhất: 2026-09-07 | Ôn gần nhất: chưa
%  Dependency: A/02, D/17. Mức độ: cốt lõi.
% ============================================================================
\documentclass[../../english.tex]{subfiles}
\begin{document}
\chapter{Bốn nhánh luyện tập (The Four Strands)}
\ptrtarget{E/19}\ptrtarget{E/19-bon-nhanh-luyen}
\dependency{\ptr{A/02} cho chín mặt của việc biết một từ; \ptr{D/17} cho việc học theo cụm.}

\begin{tomtat}
Chương này cho khung phân bổ thời gian mà mọi thứ còn lại trong phần E dựa vào. Khung đến từ Nation: một chương trình học ngôn ngữ cân bằng chia đều thời gian cho bốn nhánh --- tiếp nhận dễ, học có chủ đích, sản xuất, và luyện lưu loát. Quan sát then chốt: phần lớn người tự học ở trình độ của bạn dồn gần hết thời gian vào một nhánh và bỏ trống ba nhánh còn lại, và đó là lý do họ dừng lại ở một mức rồi không tiến nữa. Nội dung gồm: bốn nhánh và mục đích riêng của từng nhánh; bài tự chẩn đoán; khoảng cách giữa nhận ra và dùng được; luyện có chủ đích áp cho ngôn ngữ; và cách gắn tiếng Anh vào công việc thật để không phải tìm thời gian riêng. Nếu chỉ nhớ một điều: mỗi nhánh làm một việc mà ba nhánh kia \emph{không} làm được thay.
\end{tomtat}

\begin{nentang}
\kn{tiếp nhận dễ}{đọc và nghe tài liệu mà bạn hiểu gần hết --- khoảng 98\% từ đã biết; mục đích là củng cố và tự động hoá.}
\kn{học có chủ đích}{nghiên cứu ngôn ngữ một cách trực tiếp: tra từ, học cụm, phân tích cấu trúc.}
\kn{nhánh sản xuất}{viết và nói --- ép bạn truy xuất chủ động chứ không chỉ nhận ra.}
\kn{luyện lưu loát}{làm việc dễ hơn khả năng hiện có, với áp lực về tốc độ; không học cái mới.}
\kn{vốn từ chủ động}{các từ và cụm bạn dùng được khi viết; luôn nhỏ hơn nhiều so với vốn từ nhận ra.}
\kn{luyện có chủ đích}{luyện tập nhắm vào điểm yếu cụ thể, có phản hồi, ở mức khó vừa quá khả năng.}
\end{nentang}

\begin{khoigoi}
\h{Bạn đọc rất nhiều tài liệu tiếng Anh mỗi ngày mà vốn từ chủ động không tăng. Vì sao?}
\h{Vì sao \emph{luyện lưu loát} --- làm những việc dễ hơn khả năng --- lại cần thiết, nghe như lãng phí thời gian?}
\h{Vì sao bạn nhận ra hàng nghìn từ khi đọc mà chỉ dùng được vài trăm khi viết?}
\h{Đọc tài liệu chuyên ngành khó thuộc nhánh nào, và nó có thay được nhánh tiếp nhận dễ không?}
\h{Bạn bận và không có giờ trống riêng cho tiếng Anh. Giải quyết thế nào mà không phải bỏ cuộc?}
\end{khoigoi}

Hãy bắt đầu bằng câu hỏi mở đầu thứ nhất, vì nó mô tả đúng tình huống của phần lớn người làm chuyên môn.

Bạn đọc bài báo tiếng Anh hằng ngày. Bạn hiểu chúng. Nhưng khi ngồi viết, bạn vẫn dùng đi dùng lại vài chục cấu trúc quen thuộc và vẫn phải tra từ. Đây không phải vì bạn đọc chưa đủ; nó là vì đọc và viết luyện hai thứ khác nhau, và đọc \emph{không} tự động chuyển thành viết.

Nation đề xuất một khung giải thích rõ chuyện này và cho luôn cách sửa: một chương trình cân bằng cần bốn nhánh, mỗi nhánh khoảng một phần tư thời gian, và mỗi nhánh làm một việc mà ba nhánh kia không làm được thay.

\section{Bốn nhánh}

\begin{center}\footnotesize
\rowcolors{2}{white}{tblalt}
\begin{longtable}{@{}L{2.9cm}L{4.3cm}L{4.0cm}L{3.6cm}@{}}
\toprule
\textbf{Nhánh} & \textbf{Việc bạn làm} & \textbf{Nó luyện cái gì} & \textbf{Điều kiện}\\
\midrule
\endfirsthead
\toprule
\textbf{Nhánh} & \textbf{Việc bạn làm} & \textbf{Nó luyện cái gì} & \textbf{Điều kiện (tiếp)}\\
\midrule
\endhead
Tiếp nhận dễ & Đọc và nghe tài liệu hiểu gần hết & Củng cố cái đã biết; gặp lại từ trong nhiều ngữ cảnh & Khoảng 98\% từ đã biết; nếu phải tra liên tục thì sai nhánh\\
Học có chủ đích & Tra từ, ghi cụm, phân tích cấu trúc, làm bài tập & Đưa cái mới vào; sửa lỗi hoá thạch & Có mục tiêu cụ thể; không đọc lan man\\
Sản xuất & Viết và nói & Chuyển từ nhận ra sang dùng được & Phải có phản hồi, dù chỉ là tự đối chiếu\\
Luyện lưu loát & Đọc nhanh, viết nhanh, nói theo chủ đề quen & Tự động hoá; giảm tải nhận thức & Không có cái gì mới; có áp lực tốc độ\\
\bottomrule
\end{longtable}
\end{center}

Cột cuối đáng đọc kỹ, vì đó là chỗ người tự học hay làm sai mà không biết. Đọc tài liệu chuyên ngành mà cứ vài dòng phải tra từ \emph{không} phải nhánh tiếp nhận dễ --- nó là nhánh học có chủ đích, và nó không cho bạn lợi ích của nhánh một. Đây là câu trả lời cho câu hỏi mở đầu thứ tư: nó không thay được, vì hai nhánh phục vụ hai mục đích khác nhau.

\begin{figure}[H]\centering
\begin{tikzpicture}[
  q/.style={draw=steel,fill=bluebg,rounded corners=3pt,inner sep=5pt,align=center,font=\small,text width=3.5cm,minimum height=1.5cm},
  qo/.style={q,draw=accent,fill=amberbg},
  lb/.style={font=\scriptsize\itshape,color=tiercol,align=center,text width=3.4cm}]
\node[q] (n1) {\textbf{1. Tiếp nhận dễ}\\ {\scriptsize đọc/nghe hiểu gần hết}};
\node[q,right=0.5cm of n1] (n2) {\textbf{2. Học có chủ đích}\\ {\scriptsize tra, ghi cụm, phân tích}};
\node[qo,below=0.5cm of n1] (n3) {\textbf{3. Sản xuất}\\ {\scriptsize viết và nói}};
\node[qo,below=0.5cm of n2] (n4) {\textbf{4. Lưu loát}\\ {\scriptsize làm dễ, làm nhanh}};
\node[lb,right=0.45cm of n2] {phần lớn người tự học ở đây};
\node[lb,right=0.45cm of n4,color=reddark] {hai nhánh thường bị bỏ trống --- và đó là chỗ khoảng cách hình thành};
\end{tikzpicture}
\caption{Bốn nhánh, mỗi nhánh khoảng một phần tư thời gian. Người làm chuyên môn thường dồn vào nhánh 1 và 2 (đọc tài liệu và tra từ) và gần như không chạm hai nhánh dưới --- đúng hai nhánh quyết định khả năng \emph{dùng}.}
\label{fig:bon-nhanh}
\end{figure}

\section{Vì sao cần luyện lưu loát}

Câu hỏi mở đầu thứ hai nghe hợp lý: nếu bạn làm những việc dễ hơn khả năng, bạn có học được gì đâu?

Câu trả lời: bạn không học cái mới, nhưng bạn \emph{làm cho cái đã biết trở nên rẻ}. Một cụm mà bạn phải nhớ ra mất hai giây chiếm mất chú ý của bạn; cùng cụm đó khi đã tự động thì không chiếm gì cả, và phần chú ý được giải phóng chuyển sang nội dung. Đây chính là cơ chế mà \ptr{D/17} mô tả từ phía người nghe, giờ nhìn từ phía người nói.

Hệ quả thực hành quan trọng: nếu bạn thấy mình \emph{biết} rất nhiều mà viết vẫn chậm và mệt, vấn đề gần như chắc chắn nằm ở nhánh này chứ không ở kiến thức. Cách chữa không phải học thêm mà là làm nhiều lần những việc bạn đã làm được, với áp lực tốc độ nhẹ: đọc một trang dễ trong thời gian định trước; viết mười lăm phút liên tục không dừng để sửa; nói lại nội dung một bài báo trong ba phút.

\section{Nhận ra và dùng được}

Câu hỏi mở đầu thứ ba chạm vào một sự thật đo được: vốn từ nhận ra của mọi người --- kể cả người bản ngữ --- lớn hơn nhiều lần vốn từ chủ động. Đây không phải khuyết điểm cần sửa hoàn toàn; nó là trạng thái bình thường. Nhưng khoảng cách đó có thể thu hẹp, và cách thu hẹp là điều đáng nói.

Cơ chế của khoảng cách: nhận ra chỉ đòi bạn nối từ hình thức tới nghĩa khi hình thức được cung cấp. Dùng được đòi bạn đi ngược --- từ ý nghĩ trong đầu tới hình thức, mà không có gợi ý nào. Hướng thứ hai khó hơn nhiều và chỉ luyện được bằng cách \emph{đi theo hướng đó}, tức là nhánh sản xuất.

Từ đó ra một nguyên tắc thực dụng cho việc ghi chép: mọi thứ bạn ghi nên được thiết kế để \emph{truy xuất theo hướng dùng}. Thẻ ghi nhớ có mặt trước là từ tiếng Anh và mặt sau là nghĩa tiếng Việt chỉ luyện hướng nhận ra --- hướng bạn vốn đã mạnh. Đảo lại, hoặc tốt hơn là dùng một câu có chỗ trống lấy từ tài liệu thật, luyện đúng hướng bạn yếu.

\begin{cannho}
Phép thử ba mươi giây: mở một trang bạn vừa đọc, gấp lại, và viết ba câu tóm tắt nội dung bằng tiếng Anh. Nếu bạn hiểu trang đó dễ dàng mà không viết nổi ba câu, bạn vừa đo được chính khoảng cách này --- và bạn biết mình cần nhánh nào. Đây cũng là một bài luyện tốt, không chỉ là bài kiểm.
\end{cannho}

\section{Luyện có chủ đích áp cho ngôn ngữ}

Nghiên cứu về việc đạt trình độ cao trong nhiều lĩnh vực chỉ ra rằng số giờ tích luỹ không đủ để giải thích kết quả; cái phân biệt là \emph{loại} luyện tập. Luyện có chủ đích có ba đặc điểm: nhắm vào một điểm yếu cụ thể, ở mức khó vừa quá khả năng hiện tại, và có phản hồi nhanh.

Áp cho ngôn ngữ, ba đặc điểm này cho ba yêu cầu cụ thể.

\emph{Nhắm điểm yếu cụ thể.} ``Luyện viết'' không phải mục tiêu; ``viết phần mô tả kết quả mà không dùng danh từ hoá thừa'' thì có. Bốn phần trước của quyển này chính là danh mục các điểm yếu có thể nhắm: mạch văn (\ptr{C/12}), mức cam kết (\ptr{C/14}), kết hợp từ (\ptr{D/17}).

\emph{Mức khó vừa quá khả năng.} Với đọc, đó là tài liệu bạn hiểu khoảng chín phần mười. Với viết, đó là thể loại bạn đã viết được nhưng chưa thạo.

\emph{Phản hồi nhanh.} Đây là chỗ khó nhất cho người tự học, vì không có ai chấm bài. Ba nguồn phản hồi khả thi: đối chiếu bản viết của bạn với một bài mẫu cùng thể loại; tra bằng ngữ liệu ngành (\ptr{D/16} mục 4); và mô hình ngôn ngữ, với các giới hạn mà \ptr{E/20} nói rõ.

\begin{cambay}
Đừng nhầm \emph{lượng} tiếp xúc với luyện tập. Nghe podcast tiếng Anh tám tiếng một tuần trong lúc làm việc khác là nhánh một ở mức rất loãng, và nó gần như không di chuyển gì cả nếu ba nhánh kia trống. Điều này giải thích một hiện tượng phổ biến: người sống trong môi trường tiếng Anh nhiều năm mà vẫn mắc đúng những lỗi cũ --- tiếp xúc nhiều, nhưng không có nhánh hai để sửa lỗi và không có phản hồi.
\end{cambay}

\section{Gắn tiếng Anh vào công việc thật}

Câu hỏi mở đầu thứ năm là câu hỏi quyết định việc bạn có duy trì được hay không, và câu trả lời không phải ``quản lý thời gian tốt hơn''.

Nguyên tắc: đừng tạo thêm việc mới; hãy \emph{đổi ngôn ngữ} của việc bạn vốn đã làm. Bạn vốn đã đọc tài liệu, viết báo cáo, ghi chú, trình bày kết quả. Mỗi hoạt động đó có thể trở thành một nhánh mà không tốn thêm giờ nào.

\begin{center}\footnotesize
\rowcolors{2}{white}{tblalt}
\begin{longtable}{@{}L{4.6cm}L{4.4cm}L{5.4cm}@{}}
\toprule
\textbf{Việc bạn vốn đã làm} & \textbf{Thành nhánh nào} & \textbf{Thay đổi nhỏ cần làm}\\
\midrule
\endfirsthead
\toprule
\textbf{Việc bạn vốn đã làm} & \textbf{Thành nhánh nào} & \textbf{Thay đổi nhỏ cần làm (tiếp)}\\
\midrule
\endhead
Đọc bài báo trong ngành & 1 hoặc 2 & Chọn bài dễ cho nhánh 1; bài khó thì ghi cụm (\ptr{D/17})\\
Ghi chú công việc & 3 & Ghi bằng tiếng Anh, kể cả khi nghĩ bằng tiếng Việt\\
Viết nhật ký kỹ thuật hằng ngày & 3 \(+\) 4 & 15 phút, không dừng để tra từ; sửa sau\\
Chuẩn bị báo cáo nhóm & 3 & Viết dàn ý tiếng Anh trước khi nói tiếng Việt\\
Xem hội nghị, khoá học & 1 & Ba câu hỏi trước khi xem (\ptr{B/09})\\
Đọc lại bài mình viết & 3 \(+\) 2 & Áp danh sách kiểm ở \ptr{C/15}\\
Giải thích công việc cho đồng nghiệp & 4 & Thử một lần bằng tiếng Anh trong ba phút\\
\bottomrule
\end{longtable}
\end{center}

Điểm mấu chốt của bảng này: không dòng nào đòi thêm giờ. Chúng đòi \emph{đổi ngôn ngữ mặc định} của một hoạt động đã có. Đó là lý do cách này duy trì được khi bạn bận, trong khi ``mỗi ngày học một giờ tiếng Anh'' thì không.

\section{Gốc rễ: từ ``học ngôn ngữ'' tới ``dùng ngôn ngữ''}

Cách dạy ngoại ngữ trong phần lớn thế kỷ 19 và đầu thế kỷ 20 tập trung vào ngữ pháp và dịch: học luật, học danh sách từ, dịch câu. Mục tiêu ngầm là \emph{đọc được văn bản cổ điển}, không phải giao tiếp --- nên phương pháp khớp với mục tiêu của thời đó, dù nó bị chê nhiều về sau.

Từ giữa thế kỷ 20, khi mục tiêu chuyển sang giao tiếp, con lắc đảo chiều mạnh. Một cực đoan có ảnh hưởng lớn là lập luận rằng ngôn ngữ được \emph{thụ đắc} qua tiếp xúc với thông tin hiểu được, chứ không được \emph{học} qua luật, và rằng việc học có ý thức đóng vai trò rất hạn chế\cite{krashen1985}. Lập luận này giải thích tốt một số hiện tượng --- đặc biệt là việc trẻ em học tiếng mẹ đẻ --- và nó đúng ở chỗ nhấn mạnh vai trò của lượng tiếp xúc, thứ mà cách dạy ngữ pháp--dịch bỏ quên.

Nhưng nó không giải thích được một quan sát bền bỉ: người lớn sống trong môi trường ngôn ngữ nhiều năm vẫn giữ nguyên một số lỗi, và họ không tự sửa dù tiếp xúc rất nhiều. Điều đó gợi ý rằng tiếp xúc là cần nhưng không đủ, và rằng chú ý có ý thức vào hình thức ngôn ngữ có vai trò thật.

Khung bốn nhánh của Nation là một tổng hợp thực dụng của cuộc tranh luận đó\cite{nation2013,nation2006}. Nó giữ lại vai trò trung tâm của tiếp xúc --- nhánh một --- nhưng không coi nó là đủ; nó dành một phần tư cho việc học có ý thức, một phần tư cho sản xuất, và một phần tư cho tự động hoá. Điểm mạnh của khung này với người tự học không nằm ở lý thuyết mà ở chỗ nó \emph{đo được}: bạn nhìn vào một tuần của mình và biết ngay nhánh nào trống.

Một dòng nghiên cứu song song đóng góp phần còn lại. Các nghiên cứu về trí nhớ chỉ ra rằng việc \emph{truy xuất} một thứ ra khỏi trí nhớ củng cố nó mạnh hơn việc đọc lại nó\cite{roediger2006}, và rằng ôn tập giãn cách hiệu quả hơn ôn dồn\cite{cepeda2006}. Hai kết quả này là cơ sở cho toàn bộ phần \emph{Ôn nhanh} ở cuối mỗi chương của quyển sách này, và cho lời khuyên ở mục 3 về hướng thiết kế thẻ ghi nhớ.

\begin{gocre}
\gr{Thế kỷ 19 --- ngữ pháp và dịch}{Mục tiêu là đọc văn bản cổ điển.}{Học luật và danh sách từ; khớp mục tiêu thời đó nhưng không tạo được khả năng dùng.}
\gr{Giữa thế kỷ 20 --- chuyển sang giao tiếp}{Người học cần dùng ngôn ngữ, không chỉ đọc.}{Nhấn mạnh lượng tiếp xúc và thông tin hiểu được\cite{krashen1985}.}
\gr{Quan sát ngược lại}{Vì sao người sống lâu trong môi trường vẫn giữ lỗi cũ?}{Tiếp xúc là cần nhưng không đủ; chú ý có ý thức vào hình thức có vai trò thật.}
\gr{Nation --- bốn nhánh}{Chia thời gian thế nào cho cân?}{Bốn nhánh, mỗi nhánh một phần tư\cite{nation2013,nation2006}; đo được và tự chẩn đoán được.}
\gr{Nghiên cứu về truy xuất}{Đọc lại hay tự nhớ lại thì nhớ lâu hơn?}{Truy xuất củng cố mạnh hơn đọc lại\cite{roediger2006} \(\Rightarrow\) phần Ôn nhanh của quyển này.}
\gr{Ôn tập giãn cách}{Ôn dồn hay ôn rải?}{Rải hiệu quả hơn dồn\cite{cepeda2006} \(\Rightarrow\) thiết kế lịch ở \ptr{E/21}.}
\end{gocre}

\section{Liên hệ ngang}

\begin{lienhe}
\lh{Thể thao}{Tập kỹ thuật, tập thể lực, thi đấu}{Cùng cấu trúc bốn nhánh: không vận động viên nào chỉ thi đấu, và cũng không ai chỉ tập kỹ thuật. Nhánh lưu loát tương ứng với việc tập ở cường độ dưới ngưỡng.}
\lh{Âm nhạc}{Luyện gam, học bài mới, biểu diễn}{Luyện gam là nhánh bốn --- không học cái mới nhưng làm cái đã biết trở nên rẻ.}
\lh{Lập trình}{Đọc mã, viết mã, làm bài luyện}{Đọc nhiều mã không tự thành viết được mã; cùng khoảng cách nhận ra--dùng được ở mục 3.}
\lh{Giáo dục}{Hiệu ứng kiểm tra}{Cơ sở cho phần Ôn nhanh: tự nhớ lại củng cố mạnh hơn đọc lại\cite{roediger2006}.}
\lh{Quản lý thói quen}{Gắn việc mới vào việc đã có}{Nguyên tắc ở mục 5: đổi ngôn ngữ của hoạt động sẵn có thay vì thêm hoạt động mới.}
\lh{Học nghề}{Học qua tham gia cộng đồng}{Bổ sung cho ẩn dụ ``thu nhận kiến thức'' mà \ptr{D/18} phê bình: bạn học tiếng Anh khoa học bằng cách làm những việc cộng đồng đó làm.}
\end{lienhe}

\begin{traloi}
\tl{Vì đọc và viết luyện hai thứ khác nhau, và đọc không tự chuyển thành viết. Đọc chỉ đòi bạn đi từ hình thức tới nghĩa khi hình thức được cung cấp sẵn; viết đòi đi ngược lại --- từ ý trong đầu tới hình thức, không có gợi ý nào. Hướng thứ hai chỉ luyện được bằng cách đi theo hướng đó, tức nhánh sản xuất. Trong khung bốn nhánh, bạn đang dồn thời gian vào nhánh 1 và 2 và bỏ trống nhánh 3.}
\tl{Vì nó không dạy cái mới nhưng làm cái đã biết trở nên \emph{rẻ}. Một cụm phải nhớ ra mất hai giây chiếm mất chú ý của bạn; cùng cụm đó khi đã tự động thì không chiếm gì, và phần chú ý được giải phóng chuyển sang nội dung. Nếu bạn biết nhiều mà viết vẫn chậm và mệt, vấn đề gần như chắc chắn ở nhánh này chứ không ở kiến thức --- và cách chữa là làm nhiều lần những việc đã làm được, có áp lực tốc độ nhẹ.}
\tl{Vì hai việc đòi hai hướng truy xuất khác nhau. Nhận ra chỉ cần nối hình thức tới nghĩa khi hình thức được cung cấp; dùng được cần đi từ ý tới hình thức mà không có gợi ý. Khoảng cách này là bình thường --- người bản ngữ cũng có --- nhưng thu hẹp được bằng nhánh sản xuất. Hệ quả cho việc ghi chép: thẻ có mặt trước là từ tiếng Anh chỉ luyện hướng bạn vốn đã mạnh; hãy đảo lại, hoặc dùng một câu có chỗ trống lấy từ tài liệu thật.}
\tl{Nó thuộc nhánh học có chủ đích, không thuộc nhánh tiếp nhận dễ --- vì điều kiện của nhánh một là bạn hiểu gần hết, khoảng 98\% từ đã biết. Nếu cứ vài dòng phải tra từ thì bạn đang ở nhánh hai. Và nó không thay được nhánh một, vì hai nhánh phục vụ hai mục đích khác nhau: nhánh hai đưa cái mới vào, nhánh một củng cố và tự động hoá cái đã có qua việc gặp lại trong nhiều ngữ cảnh.}
\tl{Đừng tạo thêm việc mới --- hãy đổi \emph{ngôn ngữ} của việc bạn vốn đã làm. Ghi chú công việc bằng tiếng Anh; viết nhật ký kỹ thuật 15 phút không dừng để tra từ; chuẩn bị dàn ý báo cáo bằng tiếng Anh; đặt ba câu hỏi trước khi xem một khoá học; áp danh sách kiểm ở \ptr{C/15} khi đọc lại bài mình. Không việc nào đòi thêm giờ. Đó là lý do cách này duy trì được khi bạn bận, còn ``mỗi ngày học một giờ'' thì không.}
\end{traloi}

\begin{dongian}
Có một chuyện làm nhiều người nản: đọc tài liệu tiếng Anh mỗi ngày suốt nhiều năm, hiểu hết, mà đến lúc viết thì vẫn chậm và vẫn dùng đi dùng lại vài chục câu quen thuộc.

Đây không phải do bạn đọc chưa đủ. Đọc và viết luyện hai thứ khác nhau. Khi đọc, bạn chỉ cần nhận ra --- từ đã nằm sẵn trước mắt bạn. Khi viết, bạn phải tự lôi nó ra từ trong đầu, không có gợi ý nào. Chiều thứ hai khó hơn nhiều, và chỉ luyện được bằng cách làm đúng chiều đó.

Có một cách chia thời gian rất gọn giúp bạn thấy mình đang thiếu gì. Bốn nhánh, mỗi nhánh khoảng một phần tư:

Nhánh một: đọc và nghe những thứ bạn hiểu gần hết. Nhánh hai: tra từ, ghi cụm, học có chủ đích. Nhánh ba: viết và nói. Nhánh bốn: làm những việc dễ hơn khả năng, nhưng nhanh.

Người làm chuyên môn gần như luôn dồn hết vào nhánh một và hai, rồi để trống hai nhánh dưới --- đúng hai nhánh quyết định việc bạn \emph{dùng} được hay không.

Nhánh bốn nghe vô lý nhất nên nói thêm. Làm việc dễ thì học được gì? Bạn không học cái mới, nhưng bạn làm cho cái đã biết trở nên \emph{không tốn sức}. Một cụm mà bạn phải nghĩ hai giây mới nhớ ra sẽ chiếm mất sự tập trung của bạn; khi nó thành phản xạ, phần tập trung đó dành cho nội dung. Nếu bạn biết nhiều mà viết vẫn mệt, đây mới là chỗ thiếu, không phải kiến thức.

Còn chuyện không có thời gian? Đừng cố tìm thêm giờ. Hãy đổi ngôn ngữ của những việc bạn vốn đã làm: ghi chú công việc bằng tiếng Anh, viết dàn ý báo cáo bằng tiếng Anh, mỗi ngày viết mười lăm phút nhật ký kỹ thuật mà không dừng lại tra từ. Không tốn thêm phút nào, và đó là lý do nó duy trì được.
\motcau{Đọc không tự thành viết --- phải luyện đúng chiều.}
\chuadongian{Tỉ lệ một phần tư cho mỗi nhánh là hướng dẫn, không phải kết quả đo; tỉ lệ tối ưu tuỳ mục tiêu và trình độ.}
\end{dongian}

\begin{onbai}
\ar[3]{Bốn nhánh}{Tiếp nhận dễ; học có chủ đích; sản xuất; luyện lưu loát --- mỗi nhánh khoảng một phần tư.}
\ar[3]{Nguyên tắc trung tâm}{Mỗi nhánh làm một việc mà ba nhánh kia không làm thay được.}
\ar[3]{Điều kiện của nhánh một}{Hiểu gần hết, khoảng 98\% từ đã biết; phải tra liên tục nghĩa là bạn đang ở nhánh hai.}
\ar[3]{Nhánh nào hay bị bỏ trống}{Sản xuất và lưu loát --- đúng hai nhánh quyết định khả năng dùng.}
\ar[3]{Vì sao cần luyện lưu loát}{Không dạy cái mới nhưng làm cái đã biết trở nên rẻ, giải phóng chú ý cho nội dung.}
\ar[2]{Khoảng cách nhận ra--dùng được}{Nhận ra đi từ hình thức tới nghĩa; dùng được đi ngược lại, không có gợi ý.}
\ar[2]{Hệ quả cho thẻ ghi nhớ}{Mặt trước là từ tiếng Anh chỉ luyện hướng bạn đã mạnh; hãy đảo hoặc dùng câu có chỗ trống.}
\ar[2]{Ba đặc điểm của luyện có chủ đích}{Nhắm điểm yếu cụ thể; khó vừa quá khả năng; có phản hồi nhanh.}
\ar[2]{Ba nguồn phản hồi cho người tự học}{Đối chiếu với bài mẫu; tra ngữ liệu ngành; mô hình ngôn ngữ (có giới hạn).}
\ar[2]{Bẫy lượng tiếp xúc}{Nghe nền tám tiếng một tuần là nhánh một rất loãng; không có nhánh hai thì lỗi không tự sửa.}
\ar[2]{Nguyên tắc duy trì}{Đừng thêm việc mới; đổi ngôn ngữ của việc vốn đã làm.}
\ar[1]{Phép thử ba mươi giây}{Đọc một trang, gấp lại, viết ba câu tóm tắt bằng tiếng Anh.}
\end{onbai}

\begin{hoidap}
\qa{Tôi nên chia đúng 25\% cho mỗi nhánh không?}{Con số một phần tư là hướng dẫn để bạn không bỏ trống nhánh nào, không phải một kết quả đo được. Trong thực tế, tỉ lệ nên lệch theo mục tiêu: nếu bạn cần viết bài trong ba tháng tới, hãy tăng nhánh ba lên; nếu bạn vừa chuyển sang một lĩnh vực mới với nhiều thuật ngữ lạ, nhánh hai sẽ chiếm nhiều hơn một thời gian. Cách dùng khung đúng là dùng nó để \emph{phát hiện nhánh bằng không}, chứ không để cân đo từng phút.}
\qa{Nhánh một đòi hiểu 98\%. Tài liệu ngành tôi khó hơn thế. Lấy đâu ra tài liệu dễ?}{Ba nguồn thực tế. Bài tổng quan và bài giới thiệu trong chính ngành bạn --- chúng dùng cùng thuật ngữ nhưng câu đơn giản hơn nhiều. Sách giáo trình đại học trong lĩnh vực bạn đã vững. Và các bài viết phổ biến khoa học chất lượng cao. Điểm chung: cùng miền khái niệm, mật độ từ mới thấp. Đừng chọn tiểu thuyết --- nó dễ về khái niệm nhưng lại đầy từ vựng đời thường mà bạn không cần.}
\qa{Viết mỗi ngày mà không ai sửa thì có ích gì không?}{Có, dù ít hơn khi có phản hồi. Nhánh ba tạo ra lợi ích ngay cả khi không có người sửa, vì nó luyện đúng hướng truy xuất mà mục 3 nói tới, và nó làm lộ ra chỗ bạn không diễn đạt được --- thông tin mà việc đọc không cho bạn. Sau đó bạn có thể tự tạo phản hồi: đánh dấu ba chỗ bạn thấy bí, rồi tra chúng như một buổi nhánh hai. Đó là vòng lặp tự học hoàn chỉnh mà không cần ai.}
\qa{Tôi luyện nói được không khi không có người để nói cùng?}{Được, và với mục tiêu của bạn thì nói không phải ưu tiên cao nhất. Nhưng nếu muốn, có hai bài tập không cần người: giải thích một khái niệm trong ngành bạn trong ba phút, ghi âm lại, rồi nghe lại và đánh dấu chỗ vấp; và tóm tắt miệng một bài báo vừa đọc. Cả hai đều là nhánh ba cộng nhánh bốn, và cái thứ hai còn có lợi cho việc đọc vì nó ép bạn dựng lại lập luận (\ptr{B/06}).}
\qa{Bao lâu thì thấy tiến bộ?}{Nếu bạn đang bỏ trống nhánh ba và bốn, thay đổi thấy được thường đến trong bốn tới sáu tuần --- nhưng nó không xuất hiện dưới dạng ``biết thêm nhiều từ'' mà dưới dạng viết nhanh hơn và ít phải dừng lại hơn. Đó là lý do \ptr{E/21} đề nghị bạn đo các chỉ số cụ thể (số từ viết trong mười lăm phút, số lần khựng trên một trang) thay vì dựa vào cảm giác --- cảm giác tiến bộ trong ngôn ngữ rất không đáng tin.}
\qa{Học tiếng Anh chung hay chỉ tiếng Anh chuyên ngành?}{Với mục tiêu của bạn --- đọc tài liệu, viết bài, hiểu đúng ý người nói --- hãy tập trung vào tiếng Anh học thuật và chuyên ngành, và coi tiếng Anh đời thường là phần phụ. Lý do đã nêu ở \ptr{A/02}: từ vựng học thuật chung có tỉ lệ lợi ích trên công sức cao nhất, và từ vựng ngành thì bạn đã có sẵn khái niệm. Ngoại lệ duy nhất đáng cân nhắc là nếu bạn cần giao tiếp xã hội trong môi trường quốc tế --- lúc đó thêm một phần nhỏ cho tiếng Anh hội thoại.}
\qa{Tôi hay bỏ cuộc sau hai tuần. Làm sao duy trì?}{Nguyên nhân thường không phải thiếu ý chí mà là thiết kế: bạn đặt một hoạt động mới cạnh một ngày đã đầy. Cách chữa là mục 5 --- đổi ngôn ngữ của việc sẵn có thay vì thêm việc. Một mẹo bổ sung: đặt mức tối thiểu rất thấp cho ngày bận (một đoạn năm câu, không phải mười lăm phút) để chuỗi không bị đứt, vì đứt chuỗi mới là thứ giết thói quen chứ không phải làm ít.}
\end{hoidap}

\begin{baitap}
\bt[1]{Ghi lại một tuần thực tế của bạn và phân loại mỗi hoạt động vào bốn nhánh}{Tuần của bạn}{Cộng thời gian từng nhánh; nhánh nào bằng không?}
\bt[1]{Làm phép thử ba mươi giây với ba trang bạn vừa đọc}{Tài liệu ngành}{Đo khoảng cách nhận ra--dùng được của chính bạn.}
\bt[2]{Tìm ba nguồn tài liệu ở mức 98\% cho nhánh một}{Tài liệu ngành}{Kiểm bằng cách đọc một trang và đếm số từ phải tra.}
\bt[2]{Viết nhật ký kỹ thuật 15 phút mỗi ngày trong hai tuần, không dừng để tra từ}{Bài viết của bạn}{Đánh dấu chỗ bí, tra sau như một buổi nhánh hai.}
\bt[2]{Thiết kế lại 20 thẻ ghi nhớ theo hướng truy xuất chủ động}{Bảng thuật ngữ của bạn}{Mặt trước là câu có chỗ trống lấy từ tài liệu thật.}
\bt[3]{Chọn 3 việc trong bảng ở mục 5 và đổi sang tiếng Anh trong bốn tuần}{Công việc của bạn}{Không thêm giờ nào; chỉ đổi ngôn ngữ mặc định.}
\bt[3]{Thiết kế một buổi luyện có chủ đích 30 phút nhắm đúng một điểm yếu của bạn}{Bài viết của bạn}{Phải có: điểm yếu cụ thể, mức khó vừa phải, và nguồn phản hồi.}
\end{baitap}

\section*{Kết chương và đọc tiếp}

Chương này cho khung để bạn tự chẩn đoán: bốn nhánh, và câu hỏi ``nhánh nào của tôi đang bằng không''. Với phần lớn người làm chuyên môn, câu trả lời là nhánh ba và bốn --- và đó chính là lý do khoảng cách giữa hiểu và dùng không tự đóng lại theo thời gian.

Nhưng để luyện có chủ đích thì cần phản hồi, và người tự học không có ai chấm bài. Đây là chỗ công cụ trở nên quan trọng --- và cũng là chỗ có một rủi ro mới mà thế hệ người học trước không gặp: các công cụ hiện nay đủ tốt để \emph{làm thay} bạn, và làm thay thì bạn không học được gì. Chương \ptr{E/20} nói về cách vạch ranh giới đó.
\ifSubfilesClassLoaded{\printbibliography[title={Nguồn}]}{}
\end{document}
